% ==============================================================================
% File: fmlm_paper.tex
% Bài báo: Flow Map Language Models: One-step Language Modeling via Continuous Denoising
% Tác giả: Chanhyuk Lee, Jaehoon Yoo, Manan Agarwal, Sheel Shah, Jerry Huang,
%          Aditi Raghunathan, Seunghoon Hong, Nicholas M. Boffi, Jinwoo Kim
%          (KAIST và CMU, arXiv:2602.16813v3, 2026)
% Nguồn số liệu: reports/2602.16813_flow_map_language_models.md
% ==============================================================================

\section{Mô hình ngôn ngữ ánh xạ dòng FMLM}

% ------------------------------------------------------------------------------
% Slide 0: Tổng quan bài báo
% ------------------------------------------------------------------------------
\begin{frame}{FMLM: Tổng quan bài báo}
    \begin{itemize}
        \item \textbf{Tên bài báo:} \textit{Flow Map Language Models: One-step Language Modeling via Continuous Denoising}
        \item \textbf{Tác giả và tổ chức:} Chanhyuk Lee, Nicholas M. Boffi, Jinwoo Kim et al., KAIST và CMU, arXiv:2602.16813
        \item \textbf{Ý tưởng cốt lõi:} Chạy dòng liên tục trên biểu diễn one-hot của token, rồi chưng cất thành ánh xạ dòng để sinh văn bản chỉ với một bước.
        \item \textbf{Đóng góp nổi bật:}
        \begin{itemize}
            \item \textbf{Phản bác định kiến:} Dòng liên tục cạnh tranh được với khuếch tán rời rạc ngay cả trên dữ liệu rời rạc.
            \item \textbf{Bộ khử nhiễu hai thời điểm $\delta_{s,t}$:} Luôn nằm trên đơn hình xác suất, nên huấn luyện ánh xạ dòng bằng KL và cross-entropy thay vì sai số bình phương.
            \item \textbf{Sinh một bước:} Gen. PPL 119,34 trên LM1B, trong khi các mô hình khuếch tán rời rạc đã chưng cất vượt 1000 ở một bước.
        \end{itemize}
    \end{itemize}
\end{frame}

% ------------------------------------------------------------------------------
% Slide 1: Bối cảnh và thách thức
% ------------------------------------------------------------------------------
\begin{frame}{1. FMLM: Lỗi phân rã nhân tử của khuếch tán rời rạc}
    \small
    \textbf{Nút thắt của mô hình tự hồi quy:} Sinh tuần tự từng token, cần $L$ lần gọi mạng cho chuỗi dài $L$.
    \vspace{0.15cm}

    \textbf{Khuếch tán rời rạc buộc phải xấp xỉ:}
    $$\tilde{p}_{t|s}(\mathbf{y}_t \mid \mathbf{y}_s) := \prod_{l=1}^{L} p^{l}_{t|s}(y_t^{l} \mid \mathbf{y}_s)$$
    \begin{itemize}
        \item Giả định các token độc lập có điều kiện, chỉ đúng khi bước nhảy $t \to s$ rất nhỏ.
        \item Ví dụ đồ chơi với hai mode ``new-york'' và ``san-diego'': ít bước sẽ sinh ra cả ``new-diego'' và ``san-york''.
        \item Hậu quả ở 1 đến 4 bước: văn bản vô nghĩa với Gen. PPL trên 1000, hoặc lặp token với entropy dưới 4.
    \end{itemize}
    \vspace{0.15cm}
    \textbf{Mệnh đề 3.5:} Trên tập hữu hạn, tồn tại phân phối đích không thể thu được từ phân phối nguồn bằng một ánh xạ xác định ở mức mẫu. Vì vậy khuếch tán rời rạc không có ánh xạ dòng ở mức mẫu.
\end{frame}

% ------------------------------------------------------------------------------
% Slide 2: Mô hình dòng FLM
% ------------------------------------------------------------------------------
\begin{frame}{2. FMLM: Mô hình dòng FLM trên biểu diễn one-hot}
    \begin{itemize}
        \item \textbf{Biểu diễn:} $f(\mathbf{y}) \in \mathbb{R}^{L \times |\mathcal{V}|}$ là one-hot từng token, giải mã bằng argmax. Không cần embedding học.
        \item \textbf{Nội suy tuyến tính} giữa nhiễu Gaussian $\mathbf{x}_0$ và dữ liệu one-hot $\mathbf{x}_1$:
        $$I_t = (1-t)\mathbf{x}_0 + t\mathbf{x}_1, \qquad \mathbf{b}_t(\mathbf{x}) = \frac{D_t(\mathbf{x}) - \mathbf{x}}{1-t}$$
        \item \textbf{Học vận tốc bằng sai số bình phương thất bại:} đích $\mathbf{x}_1 - \mathbf{x}_0$ có hạng đầy đủ, mạng bị underfit khi $|\mathcal{V}| \gg d$.
        \item \textbf{Học bộ khử nhiễu thay thế:} $D_t(\mathbf{x}) = \mathbb{E}[\mathbf{x}_1 \mid I_t = \mathbf{x}]$ chính là phân phối hậu nghiệm của từng token, nằm trên đơn hình xác suất (Bổ đề 3.1).
        \item \textbf{Huấn luyện bằng cross-entropy} với đầu ra softmax theo token. Mệnh đề 3.2: nghiệm tối ưu duy nhất là $D_t$, và sai số nhỏ kéo theo khoảng cách Wasserstein $W_2$ nhỏ.
    \end{itemize}
\end{frame}

% ------------------------------------------------------------------------------
% Slide 3: Tái tham số hóa thời gian
% ------------------------------------------------------------------------------
\begin{frame}{3. FMLM: Tái tham số hóa thời gian theo tỉ lệ lỗi giải mã}
    \begin{itemize}
        \item \textbf{Vấn đề:} Tỉ lệ lỗi giải mã $P_e(t)$ giảm rất gấp gần $t = 1$ khi $|\mathcal{V}|$ lớn. Lấy $t$ đều làm lãng phí tín hiệu huấn luyện và lãng phí bước lấy mẫu.
        $$P_e(t) = \frac{1}{L}\sum_{l=1}^{L} \mathbb{P}\big(g^l(\mathbf{x}_t) \neq g^l(\mathbf{x}_1)\big)$$
        \item \textbf{Đề xuất:} Huấn luyện và lấy mẫu đều theo $\tau(t)$, sao cho mỗi bước giảm cùng một lượng lỗi giải mã:
        $$\tau(t) = 1 - \frac{|\mathcal{V}|}{|\mathcal{V}|-1}\,P_e(t)$$
        \item \textbf{Cài đặt:} Bảng tra 1000 điểm tính bằng cầu phương Gauss--Hermite, nội suy spline bậc ba. Chi phí gần như bằng không.
        \item \textbf{Hiệu quả:} Cho phép mở rộng tới $|\mathcal{V}| \approx 50\,000$. Trên LM1B, Gen. PPL của FLM: không tái tham số hóa 149,18; thời gian entropy học được 130,42; theo tỉ lệ lỗi giải mã \textbf{106,98}.
    \end{itemize}
\end{frame}

% ------------------------------------------------------------------------------
% Slide 4: Ánh xạ dòng và bộ khử nhiễu hai thời điểm
% ------------------------------------------------------------------------------
\begin{frame}{4. FMLM: Bộ khử nhiễu hai thời điểm trên đơn hình xác suất}
    \textbf{Ánh xạ dòng:} $X_{s,t}(\mathbf{x}) = \mathbf{x} + (t-s)\,\mathbf{v}_{s,t}(\mathbf{x})$, nhảy thẳng giữa hai thời điểm bất kỳ. Sinh một bước: $\hat{\mathbf{x}}_1 = X_{0,1}(\mathbf{x}_0)$.
    \begin{itemize}
        \item \textit{Khó khăn:} $X_{s,t}$ rời khỏi đơn hình do nhiễu Gaussian, nên khó huấn luyện bằng cross-entropy.
    \end{itemize}
    \vspace{0.1cm}
    \textbf{Định nghĩa:} $\delta_{s,t}(\mathbf{x}) := \mathbf{x} + (1-s)\,\mathbf{v}_{s,t}(\mathbf{x})$. Mệnh đề 3.3:
    \begin{enumerate}
        \item \textbf{Khôi phục ánh xạ dòng:} $X_{s,t}(\mathbf{x}) = \frac{1-t}{1-s}\mathbf{x} + \frac{t-s}{1-s}\,\delta_{s,t}(\mathbf{x})$.
        \item \textbf{Nằm trên đơn hình:} $\delta_{s,t}(\mathbf{x})^l \in \Delta^{|\mathcal{V}|-1}$, nên dùng được softmax.
        \item \textbf{Trùng bộ khử nhiễu trên đường chéo:} $\delta_{s,s} = D_s$.
        \item \textbf{Điều kiện nửa nhóm dạng tổ hợp lồi:}
        $\delta_{s,t} = \gamma\,\delta_{s,u} + (1-\gamma)\,\delta_{u,t}(X_{s,u})$, với $\gamma = \frac{(1-t)(u-s)}{(1-u)(t-s)} \in [0,1]$.
    \end{enumerate}
\end{frame}

% ------------------------------------------------------------------------------
% Slide 5: Chưng cất nửa nhóm
% ------------------------------------------------------------------------------
\begin{frame}{5. FMLM: Chưng cất nửa nhóm và sinh một bước}
    \textbf{Hàm mục tiêu KL nửa nhóm} với giáo viên FLM đóng băng $\hat{D}$ (Mệnh đề 3.4):
    \begin{small}
    $$\mathcal{L}_{\text{KL}} = \mathbb{E}\!\left[\sum_l \mathrm{KL}\big(\bar\delta^{l}_{s,t} \,\|\, \hat\delta^{l}_{s,t}(I_s)\big)\right] + \mathbb{E}\!\left[\sum_l \mathrm{KL}\big(\hat D^{l}_t(I_t) \,\|\, \hat\delta^{l}_{t,t}(I_t)\big)\right]$$
    \end{small}
    với mục tiêu giáo viên $\bar\delta_{s,t} = \mathrm{sg}\big(\gamma\,\hat\delta_{s,u}(I_s) + (1-\gamma)\,\hat\delta_{u,t}(\hat X_{s,u})\big)$.
    \begin{itemize}
        \item \textbf{Lấy mẫu bộ ba $(s,u,t)$:} chọn độ dài bước $h$, điểm đầu $\tau(s)$, điểm cuối $\tau(t) = \tau(s)+h$, điểm giữa $u$ là trung điểm. Cách này tốt nhất trong các cách thử.
        \item \textbf{Lấy mẫu biên:} cặp $(s,t) = (0,1)$ với xác suất $1/32$, vì $\tau$ phẳng gần $t=0$. Bỏ kỹ thuật này, Gen. PPL một bước tăng từ 119,34 lên 142,61.
        \item \textbf{Cấu hình:} chưng cất 100 nghìn bước, chia đôi batch cho đường chéo và ngoài đường chéo. Lấy mẫu dùng thuật toán $\gamma$ của Kim et al.
    \end{itemize}
\end{frame}

% ------------------------------------------------------------------------------
% Slide 6: Kết quả FLM nhiều bước
% ------------------------------------------------------------------------------
\begin{frame}{6. FMLM: Kết quả của FLM ở 1024 bước}
    \textbf{Gen. PPL theo GPT-2 Large, 1024 bước. Entropy của dữ liệu: LM1B 4,31; OWT 5,44.}
    \begin{table}[h]
        \centering
        \small
        \begin{tabular}{lcccc}
            \toprule
            \textbf{Mô hình} & \textbf{LM1B PPL} & \textbf{LM1B Ent} & \textbf{OWT PPL} & \textbf{OWT Ent} \\
            \midrule
            RDLM & 268,21 & 4,33 & -- & -- \\
            CANDI & 120,99 & 4,35 & 143,13 & 5,71 \\
            MDLM & 109,21 & 4,32 & 105,15 & 5,63 \\
            Duo & 98,14 & 4,31 & 77,69 & 5,55 \\
            \midrule
            \textbf{FLM} & \textbf{96,91} & 4,29 & \textbf{62,23} & 5,33 \\
            \bottomrule
        \end{tabular}
    \end{table}
    \vspace{-0.2cm}
    \begin{itemize}
        \item Trên LM1B, FLM chỉ hơn Duo khoảng 1,3\%. Trên OWT, FLM hơn rõ nhưng entropy thấp hơn dữ liệu, tác giả tự thừa nhận có đánh đổi nhỏ.
        \item Trên OWT, các đối chứng dùng checkpoint công bố chính thức, không huấn luyện lại cùng cấu hình.
        \item Mở rộng cỡ mô hình trên OWT, 150 nghìn bước: 179M, 424M, 870M cho Gen. PPL 75,57, 65,94, 61,59.
    \end{itemize}
\end{frame}

% ------------------------------------------------------------------------------
% Slide 7: Kết quả sinh ít bước
% ------------------------------------------------------------------------------
\begin{frame}{7. FMLM: Kết quả sinh một đến bốn bước}
    \small
    \textbf{Gen. PPL, thấp hơn là tốt hơn. Dấu sao đỏ: entropy dưới 4,0, tức lặp token.}
    \begin{table}[h]
        \centering
        \footnotesize
        \begin{tabular}{lccccc}
            \toprule
            \textbf{Tập, số bước} & \textbf{MDLM+SDTT} & \textbf{MDLM+Di4C} & \textbf{Duo+DCD} & \textbf{Duo+Di4C} & \textbf{FMLM} \\
            \midrule
            LM1B, 1 bước & 1429,48 & 1217,10 & 1224,52 & 292,94\textcolor{red}{$^*$} & \textbf{119,34} \\
            LM1B, 2 bước & 602,14 & 621,59 & 520,08 & 247,69\textcolor{red}{$^*$} & \textbf{110,19} \\
            LM1B, 4 bước & 241,01 & 247,32 & 210,88 & 150,67 & \textbf{98,76} \\
            \midrule
            OWT, 1 bước & 1260,86 & 1298,80 & 5743,29 & 370,51\textcolor{red}{$^*$} & \textbf{168,30} \\
            OWT, 2 bước & 877,22 & 758,23 & 891,16 & 210,22 & \textbf{133,29} \\
            OWT, 4 bước & 339,73 & 239,27 & 250,86 & 154,67 & \textbf{111,31} \\
            \bottomrule
        \end{tabular}
    \end{table}
    \vspace{-0.2cm}
    \begin{itemize}
        \item Entropy của FMLM một bước: LM1B 4,16 và OWT 5,17, thấp hơn dữ liệu thật (4,31 và 5,44).
        \item Một bước của FMLM tương đương các mô hình đối chứng ở 8 đến 16 bước trên LM1B và 4 đến 8 bước trên OWT. Hình 1 báo tăng tốc 8,3 lần.
    \end{itemize}
\end{frame}

% ------------------------------------------------------------------------------
% Slide 8: Dẫn hướng
% ------------------------------------------------------------------------------
\begin{frame}{8. FMLM: Dẫn hướng lúc suy luận}
    \small
    \textbf{Tự dẫn hướng} ngoại suy vận tốc ra xa mô hình yếu, mô hình yếu là bản có dropout 0,1:
    $$\hat{\mathbf{b}}_t^{\text{guided}} = \hat{\mathbf{b}}_t^{\text{weak}} + \eta\big(\hat{\mathbf{b}}_t - \hat{\mathbf{b}}_t^{\text{weak}}\big)$$
    \begin{itemize}
        \item FLM ổn định đến $\eta = 100$. Tại $\eta = 50$, Gen. PPL giảm từ 96,91 xuống \textbf{51,62}, nhưng entropy giảm còn 4,00.
        \item Duo và MDLM sụp đổ từ $\eta \ge 10$.
    \end{itemize}
    \vspace{0.1cm}
    \textbf{Dẫn hướng theo phần thưởng} dựa trên nhìn trước điểm cuối của ánh xạ dòng:
    $$\mathbf{x}_{t_{n+1}} = X_{t_n,t_{n+1}}(\mathbf{x}_{t_n}) + \lambda\,\nabla_{\mathbf{x}_{t_n}}\, r\big(X_{t_n,1}(\mathbf{x}_{t_n})\big)$$
    \begin{itemize}
        \item Bộ phân loại phần thưởng chỉ cần huấn luyện trên văn bản sạch. Khuếch tán rời rạc phải huấn luyện trên dữ liệu bị nhiễu ở mọi mức.
        \item Ở 2 bước, chủ đề: FMLM đạt PPL 131,8 và phần thưởng 0,670; MDLM+SDTT đạt 888,8 và 0,274. Tiêu chí ngữ pháp: phần thưởng của mọi phương pháp đều thấp.
    \end{itemize}
\end{frame}

% ------------------------------------------------------------------------------
% Slide 9: Sudoku và đa dạng
% ------------------------------------------------------------------------------
\begin{frame}{9. FMLM: Ràng buộc logic trên Sudoku và độ đa dạng}
    \textbf{Tỉ lệ lời giải hợp lệ, phần trăm, ở 1, 2, 4 và 1024 bước:}
    \begin{table}[h]
        \centering
        \footnotesize
        \begin{tabular}{lcc}
            \toprule
            & \textbf{Sinh không điều kiện} & \textbf{Giải puzzle, tập kiểm tra} \\
            \midrule
            MDLM+SDTT & 0 / 0 / 0,02 / 92,19 & 0 / 2,93 / 8,40 / 98,63 \\
            Duo+DCD & 0 / 0,39 / 17,19 / 91,41 & 0 / 30,85 / 73,82 / 96,88 \\
            \textbf{FMLM} & \textbf{9,38 / 57,03 / 79,69 / 94,92} & \textbf{39,84 / 85,15 / 92,19 / 98,83} \\
            \bottomrule
        \end{tabular}
    \end{table}
    \vspace{-0.2cm}
    \begin{itemize}
        \item Mọi lưới hợp lệ đều duy nhất và không trùng tập huấn luyện. Ở 1024 bước, các đối chứng ngang FMLM, nên lợi thế nằm ở vùng ít bước.
        \item \textbf{Độ đa dạng một bước:} Self-BLEU của FMLM là 0,073 trên LM1B và 0,121 trên OWT, dữ liệu thật là 0,047 và 0,046. Trọng tài GPT-4.1 cho FMLM thắng 39 đến 42\% về đa dạng. Kết luận: không sụp đổ mode, nhưng kém đa dạng hơn dữ liệu thật, rõ nhất trên OWT.
    \end{itemize}
\end{frame}

% ------------------------------------------------------------------------------
% Slide 10: Nghiên cứu loại bỏ thành phần
% ------------------------------------------------------------------------------
\begin{frame}{10. FMLM: Nghiên cứu loại bỏ thành phần}
    \textbf{Mô hình dòng FLM, Gen. PPL trên LM1B ở 1024 bước:}
    \begin{itemize}
        \item \textbf{Hàm mục tiêu:} vận tốc với sai số bình phương 3801,36; bộ khử nhiễu với sai số bình phương 129,04; thêm softmax 120,16; thêm cross-entropy \textbf{96,91}.
        \item \textbf{Biểu diễn:} one-hot 130,42, tốt hơn embedding học được (324,66), BERT-large đóng băng (262,92), embedding ngẫu nhiên (400,17).
        \item \textbf{Khung khuếch tán:} Euclid 96,91. Đơn hình cho PPL 85,07 nhưng entropy sụp còn 3,76. Riemann cho 268,21.
    \end{itemize}
    \vspace{0.1cm}
    \textbf{Ánh xạ dòng FMLM, Gen. PPL một bước:}
    \begin{itemize}
        \item \textbf{Tham số hóa:} vận tốc trung bình với sai số bình phương 199,01; bộ khử nhiễu hai thời điểm với softmax và cross-entropy \textbf{119,34}.
        \item \textbf{Mục tiêu chưng cất:} nửa nhóm 119,34; Eulerian 149,13; Lagrangian 193,08. Chưng cất tốt hơn tự chưng cất (159,53).
        \item \textbf{Nhiễu đầu vào:} Gaussian 105,48 so với đều trên đơn hình 660,33.
    \end{itemize}
\end{frame}

% ------------------------------------------------------------------------------
% Slide 11: Hạn chế, phản biện và kết luận
% ------------------------------------------------------------------------------
\begin{frame}{11. FMLM: Hạn chế, phản biện và kết luận}
    \small
    \textbf{Điểm cần thận trọng khi đọc kết quả:}
    \begin{itemize}
        \item Chỉ số chính là Gen. PPL theo GPT-2 Large cộng entropy, không có đánh giá của người.
        \item Entropy của FMLM luôn thấp hơn dữ liệu thật: một phần lợi thế PPL có thể đến từ việc dồn mẫu về mode phổ biến.
        \item Quy mô nhỏ: 179M tham số. Chi phí huấn luyện hai giai đoạn không được tính vào so sánh. One-hot tốn thêm khoảng 30\% thời gian và bộ nhớ.
    \end{itemize}
    \vspace{0.1cm}
    \textbf{Kết luận:}
    \begin{enumerate}
        \item Dòng liên tục trên one-hot cạnh tranh được với khuếch tán rời rạc ở nhiều bước và vượt rõ ở ít bước.
        \item Bộ khử nhiễu hai thời điểm đưa ánh xạ dòng lên đơn hình, nên chưng cất bằng cross-entropy và KL ổn định.
        \item Mở đường cho dẫn hướng và học tăng cường. Chưa có bằng chứng ở quy mô LLM.
    \end{enumerate}
\end{frame}
